\begin{tabular}{p{6.2cm}rlp{5.0cm}}
\toprule
Check & Value & Unit & Comment \\
\midrule
Tank diameter for H = T (0.50 m\textsuperscript{3}) & 0.86 & m & geometry implied by V; no tank drawing exists \\
Wetted wall + base area at H = T & 2.91 & m\textsuperscript{2} & A = 2.0 m\textsuperscript{2} is of the same order \\
Impeller-to-tank diameter ratio $D/T$ & 0.35 & - & within the commonly used 0.3--0.5 range \\
Impeller Reynolds number at 120 rpm (water-like) & 225,000 & - & turbulent: constant $N_P$ applicable \\
Shaft power per volume at recommended speed & 2.89 & kW/m\textsuperscript{3} & high-intensity duty; check with vendor \\
Loop pressure drop at 0.50 kg/s & 3.12 & bar & extrapolated from assumed 0.5 bar at 0.2 kg/s \\
Pump flow turndown modelled (0.05--0.50 kg/s) & 10.00 & :1 & model assumes variable-speed pump, constant efficiency \\
Jacket gap implied by $M_J$ over A & 5.00 & mm & thin; see isolated-inventory study \\
Plain annular jacket Re at 0.50 kg/s, $Re = 2\dot m/(\pi T\mu)$ & 462.5 & - & laminar: assumed U needs a baffled/spiral jacket \\
Coolant outlet maximum (baseline peak $T_J$) & 318.1 & K & about 45~$^{\circ}$C; confirm against site cooling-water limits \\
Liquid density range 300--370 K (water) & 960.0 & kg/m\textsuperscript{3} & about 996 to 960 kg/m\textsuperscript{3}; constant 1000 used ($\leq$ 4 \%) \\
Liquid Cp range 300--370 K (water) & 4.21 & kJ/(kg K) & about 4.18 to 4.21; constant 4.18 used ($<$ 1 \%) \\
\bottomrule
\end{tabular}
